We ask whether temperature scaling (TS) fitted on in-distribution validation data stays calibrated under distribution shift, on scikit-learn digits with rotation (0--60$^\circ$) and Gaussian-noise shift of increasing intensity. We reimplement and compare a single MLP, TS, MC dropout and a 5-member deep ensemble over 11 shift conditions and 5 seeds, reporting accuracy, ECE, NLL and Brier score. In distribution all systems are well calibrated and the TS gain is not statistically supported with 5 seeds. Under shift, ID-fitted TS does not remain calibrated: its ECE grows about 50-fold at 60$^\circ$ rotation and about 39-fold at noise $\sigma=1$, although it still lowers NLL relative to the uncalibrated MLP at every shifted condition. An oracle ablation fitting the temperature on shift-matched validation data brings ECE back to at most 0.092, but at strong shift only by pushing predictions towards uniform (NLL close to $\ln 10$ at 60$^\circ$). Against our pre-registered expectation, MC dropout had lower mean ECE, NLL and Brier than the deep ensemble in all 11 conditions (paired differences significant at $\alpha=0.05$, uncorrected, for NLL in all 11 conditions, ECE in 9 and Brier in 7; in distribution the ECE difference is within seed noise); an exploratory decomposition suggests this is partly an effect of dropout regularisation during training. The study is small (one dataset, one MLP, 5 seeds, untuned reimplemented baselines) and the conclusions are limited to that scale.
